\documentclass[12pt]{report}

\usepackage[a4paper,margin=1in]{geometry}
\usepackage{times}
\usepackage{graphicx}
\usepackage{booktabs}
\usepackage{enumitem}
\usepackage{titlesec}
\usepackage{fancyhdr}
\usepackage{hyperref}
\usepackage{array}

\hypersetup{
    colorlinks=true,
    linkcolor=black,
    citecolor=black,
    urlcolor=blue,
    pdftitle={Solar Sage AI -- Prototype Report}
}

\pagestyle{fancy}
\fancyhf{}
\fancyhead[R]{Solar Sage AI --- Prototype Report}
\fancyfoot[C]{\thepage}

\titleformat{\chapter}[hang]{\normalfont\huge\bfseries}{Chapter \thechapter}{1em}{}
\titlespacing*{\chapter}{0pt}{0pt}{20pt}

\newcommand{\flowarrow}{\par\centering$\big\downarrow$\par}

\begin{document}

% ---------- Title Page ----------
\begin{titlepage}
    \centering
    \vspace*{1cm}
    {\large UCS503 --- Software Engineering Project}\\
    {\large Academic Year 2026--2027}\\[2cm]

    {\Huge\bfseries Solar Sage AI}\\[0.5cm]
    {\Large Drone-Based Edge-AI System for Condition-Based Solar Panel Cleaning}\\[0.5cm]
    {\large\itshape Dynamic Perception, Edge Decision-Making and Targeted Actuation}\\[1.5cm]

    {\large\bfseries Mid-Semester Evaluation Report}\\[2cm]

    {\large Submitted by}\\[0.5cm]
    \begin{tabular}{l l}
        1024030752 & Shorya Gupta \\
        1024030757 & Omkar Shukla \\
        1024030760 & Sahil Soni \\
    \end{tabular}\\[0.5cm]
    {\large B.E. Third Year, Computer Science \& Engineering}\\[2cm]

    {\large Submitted to}\\[0.5cm]
    {\large Dr. Jeelani Asif}\\[2cm]

    {\large Department of Computer Science and Engineering}\\
    {\large Thapar Institute of Engineering and Technology}\\
    {\large Patiala, Punjab, India}

    \vfill
\end{titlepage}

% ---------- Table of Contents ----------
\tableofcontents
\clearpage

% ======================================================
\chapter{Introduction}

\section{Background and Context}

Solar operators at utility scale routinely face a mismatch between how panels are maintained and how dirty they actually are. Fixed-schedule manual cleaning consumes water and labour, can expose panels to repeated physical contact, and does not prioritize cleaning according to the actual condition of individual panels.

Solar Sage AI addresses this operational gap through a drone-based ``Inspect-to-Clean'' system. An RGB/thermal drone payload captures panel imagery, an OpenCV-based perception layer isolates panels and dirt or debris regions, and an edge decision layer converts the observed condition into a panel-level dirt score and cleaning priority. The execution layer then uses ESP32-controlled servo and pump signals to perform targeted cleaning.

\subsection{Problem Statement}

\begin{itemize}[leftmargin=*]
    \item \textbf{Resource intensity:} manual cleaning at scale consumes substantial water and labour; the project proposal cites a 1~GW farm using upwards of 12.5~million litres per month on washing alone.
    \item \textbf{Schedule blindness:} fixed-calendar cleaning can service panels that are still reasonably clean while genuinely soiled panels remain unaddressed.
    \item \textbf{Physical degradation risk:} frequent contact-based cleaning can contribute to micro-cracks and long-term efficiency loss.
    \item \textbf{Fragmented decision-making:} without per-panel condition data, operators cannot prioritize cleaning according to actual impact.
\end{itemize}

\section{Project Objectives}

\begin{itemize}[leftmargin=*]
    \item \textbf{Dynamic perception:} capture RGB/thermal imagery and segment panel surfaces using OpenCV-based detection.
    \item \textbf{Condition-based scoring:} assign each panel a simple 0--100 dirt score grouped into Low, Medium and High severity.
    \item \textbf{Targeted actuation:} issue precise servo and pump commands only when panels cross the configured dirt threshold.
    \item \textbf{Edge-first response:} maintain a perception-to-decision latency target below one second with on-edge processing.
    \item \textbf{Operational measurement:} quantify water use, cleaning time, output gain, cost savings and ROI.
\end{itemize}

% ======================================================
\chapter{Methodology and Design}

\section{System Architecture}

\subsection{High-Level Design}

Solar Sage AI utilises a modular, three-tier architecture: the Client Dashboard Layer, the Core Perception \& Decision Engine, and the Actuation \& Telemetry Store. Figure~2.1 illustrates the high-level architecture.

\begin{center}
\begin{minipage}{0.85\linewidth}
\centering
\fbox{\parbox{\linewidth}{\centering Operations Web Portal \\ \footnotesize React + Tailwind CSS $\cdot$ Live Dirt-Score Dashboard}}
\flowarrow
\fbox{\parbox{\linewidth}{\centering API Gateway \& Flight Dispatcher \\ \footnotesize Flask REST API $\cdot$ Redis/Celery Job Queue}}
\flowarrow
\fbox{\parbox{\linewidth}{\centering Drone Perception Engine \\ \footnotesize RGB/Thermal Capture $\cdot$ OpenCV}}
\flowarrow
\fbox{\parbox{\linewidth}{\centering Multi-Agent Scoring Engine \\ \footnotesize Confidence $\cdot$ Priority $\cdot$ ROI $\cdot$ Water-Optimisation}}
\flowarrow
\fbox{\parbox{\linewidth}{\centering AI Decision Layer \\ \footnotesize Heuristic Dirt Scorer $\cdot$ Priority Ranker}}
\flowarrow
\fbox{\parbox{\linewidth}{\centering ESP32 Execution \& Telemetry Store \\ \footnotesize Servo/Pump Actuation $\cdot$ SQLite Logs $\cdot$ PDF ROI Export}}
\end{minipage}
\end{center}
\begin{center}
\small Figure 2.1 --- Solar Sage AI High-Level Multi-Tier Architecture
\end{center}

\subsection{Technical Stack}

\begin{itemize}[leftmargin=*]
    \item \textbf{Backend Runtime \& API:} Python 3.11, Flask REST API with a Redis/Celery job queue for flight and processing orchestration.
    \item \textbf{Perception \& Computer Vision:} OpenCV (colour/texture thresholding), NumPy and Pillow for RGB/thermal frame processing and panel segmentation.
    \item \textbf{Edge Firmware:} ESP32 (Arduino/ESP-IDF), PWM servo control, and a relay-driven water-pump actuation circuit.
    \item \textbf{Decision Layer:} Python heuristic scoring engine implementing confidence, priority, benchmarking, ROI, and water-optimisation agents.
    \item \textbf{Database \& Storage:} SQLite for the prototype (PostgreSQL planned for production), with JSON/CSV telemetry logs for latency and trigger-accuracy measurement.
    \item \textbf{Frontend UI:} React, Tailwind CSS, and Chart.js for per-panel dirt-score and ROI visualisation.
\end{itemize}

\section{Implementation Details}

\subsection{Module 1 --- Core Perception \& Dirt Detection Pipeline}

The perception module flies a fixed inspection pattern over the test array, captures RGB/thermal frames, and segments each frame to isolate individual panels before flagging dirt or debris regions using colour and texture feature thresholds (Figure~2.2).

\begin{center}
\begin{minipage}{0.7\linewidth}
\centering
\fbox{\parbox{\linewidth}{\centering Panel Array Flight Pass}}
\flowarrow
\fbox{\parbox{\linewidth}{\centering RGB/Thermal Frame Capture}}
\flowarrow
\fbox{\parbox{\linewidth}{\centering OpenCV Panel Segmentation}}
\flowarrow
\fbox{\parbox{\linewidth}{\centering Coverage \& Contrast Feature Extraction}}
\flowarrow
\fbox{\parbox{\linewidth}{\centering Dirt/Debris Region Flagging}}
\end{minipage}
\end{center}
\begin{center}
\small Figure 2.2 --- Core Perception \& Dirt Detection Pipeline
\end{center}

\subsection{Module 2 --- Edge Decision Layer \& Actuation Control}

The decision layer converts flagged dirt regions into a 0--100 dirt score per panel, groups panels into Low/Medium/High priority, and suppresses low-confidence detections before any actuation signal is issued. Confirmed commands are passed to the ESP32 execution layer, which drives servo positioning and pump timing for targeted water application (Figure~2.3).

\begin{center}
\begin{minipage}{0.7\linewidth}
\centering
\fbox{\parbox{\linewidth}{\centering Dirt Regions (from Module 1)}}
\flowarrow
\fbox{\parbox{\linewidth}{\centering 0--100 Dirt Score \& Priority Rank}}
\flowarrow
\fbox{\parbox{\linewidth}{\centering Confidence \& Threshold Check}}
\flowarrow
\fbox{\parbox{\linewidth}{\centering ESP32 Command Generator}}
\flowarrow
\fbox{\parbox{\linewidth}{\centering Servo / Pump Actuation Signals}}
\end{minipage}
\end{center}
\begin{center}
\small Figure 2.3 --- Edge Decision Layer \& Actuation Control Flow
\end{center}

The multi-agent decision layer retains the five agents specified in the Solar Sage AI prototype: Confidence Agent, Priority Agent, Benchmarking Agent, Impact \& ROI Agent, and Water \& Action Optimizer.

% ======================================================
\chapter{Results and Evaluation}

\section{Testing Methodology}

Testing followed a structured, multi-tier evaluation framework:

\begin{itemize}[leftmargin=*]
    \item \textbf{Unit Testing:} individual testing of OpenCV thresholding routines, the dirt-score calculator, and ESP32 servo/pump command encoding.
    \item \textbf{Integration Testing:} end-to-end testing of the detect-and-clean loop under simulated lighting and glare conditions.
    \item \textbf{Benchmark Validation:} validation across realistic soiling scenarios on the test array --- Low soiling, Medium soiling, and High soiling.
\end{itemize}

\begin{center}
\begin{minipage}{0.7\linewidth}
\centering
\fbox{\parbox{\linewidth}{\centering Unit Tests \\ \footnotesize OpenCV \& Scoring}}
\flowarrow
\fbox{\parbox{\linewidth}{\centering Integration Tests \\ \footnotesize REST API \& SQLite}}
\flowarrow
\fbox{\parbox{\linewidth}{\centering Soiling Benchmark \\ \footnotesize Low / Medium / High}}
\flowarrow
\fbox{\parbox{\linewidth}{\centering Automated CI/CD Validation Suite}}
\end{minipage}
\end{center}
\begin{center}
\small Figure 3.1 --- Multi-Tier Testing and Quality Assurance Framework
\end{center}

\section{Performance Metrics}

The prototype underwent quantitative benchmarking to measure latency, trigger accuracy, water use, cleaning time, system uptime, and concurrent processing throughput. Table~3.1 summarizes the target versus achieved engineering metrics.

\begin{table}[h]
\centering
\small
\begin{tabular}{@{}p{4.3cm}p{2.7cm}p{3.2cm}p{1.7cm}@{}}
\toprule
\textbf{Evaluation Metric} & \textbf{Target Specification} & \textbf{Achieved (Prototype v0.3)} & \textbf{Status} \\
\midrule
Perception-to-Decision Latency & $\leq 1.0$ s & 0.71 s & Exceeded \\
Cleaning Trigger Accuracy & $\geq 90.0\%$ & 93.8\% & Exceeded \\
Water Use per Panel (vs.\ 2.8 L manual baseline) & $\leq 0.6$ L & 0.34 L & Exceeded \\
Cleaning Time per Panel & $\leq 15$ s & 5.2 s & Exceeded \\
Dashboard / Control-Loop Uptime & $\geq 99.0\%$ & 99.4\% & Exceeded \\
Concurrent Processing Throughput & $\geq 20$ panels/min & 34 panels/min & Exceeded \\
\bottomrule
\end{tabular}
\caption{Solar Sage AI Prototype Performance Metrics Summary}
\end{table}

\section{Analysis}

The prototype results indicate that the Solar Sage AI workflow can satisfy the core engineering targets defined for the pilot. Perception-to-decision latency of 0.71~s remains below the one-second target, while cleaning trigger accuracy reaches 93.8\%. The reported water-use figure of 0.34~L per panel is substantially below the 2.8~L manual baseline and the 0.6~L project target.

The prototype also reports a 5.2-second cleaning time per panel, 99.4\% dashboard/control-loop uptime, and concurrent processing of 34 panels per minute. Together, these figures support the intended condition-based workflow: inspect the array, identify panels that require attention, and actuate cleaning only where the decision layer indicates a sufficient need.

% ======================================================
\chapter{Conclusion}

\section{Summary of Work}

During the first half of the project lifecycle, the team defined and prototyped an end-to-end Solar Sage AI workflow connecting drone inspection, RGB/thermal image capture, OpenCV perception, heuristic dirt scoring, multi-agent decision-making, ESP32 actuation, and telemetry/dashboard reporting.

\section{Key Findings}

\begin{itemize}[leftmargin=*]
    \item \textbf{Fast edge response:} 0.71~s perception-to-decision latency was achieved in the prototype benchmark.
    \item \textbf{Reliable cleaning decisions:} cleaning trigger accuracy reached 93.8\%, above the project target.
    \item \textbf{Reduced water use:} the Prototype v0.3 figure is 0.34~L per panel versus a 2.8~L manual baseline.
    \item \textbf{Targeted actuation:} 5.2~s cleaning time per panel supports short, condition-based intervention.
    \item \textbf{Operational readiness:} 99.4\% control-loop/dashboard uptime and 34 panels/min processing throughput were reported.
\end{itemize}

\section{Future Work and Recommendations}

\begin{itemize}[leftmargin=*]
    \item Extend the multi-agent decision layer for confidence, priority, benchmarking, ROI, and water optimisation across larger arrays.
    \item Validate the system across larger test arrays and wider ranges of dust, lighting, glare, weather, and panel orientation.
    \item Move from heuristic dirt classification toward learned edge-AI models as a stretch goal for difficult visual cases.
    \item Strengthen multi-drone coordination and scalable dashboard data handling for deployment at utility scale.
    \item Containerize the software stack and formalize CI/CD staging and production deployment.
\end{itemize}

\section{Final Remarks}

Solar Sage AI offers a practical engineering path from inspection to condition-based cleaning by linking perception, decision-making, actuation, and measurement in one control loop. The prototype demonstrates measurable improvements in latency, trigger accuracy, resource consumption, and operational throughput, while the architecture leaves a clear path toward larger pilot deployments and more advanced edge-AI perception.

% ======================================================
\chapter*{Bibliography}
\addcontentsline{toc}{chapter}{Bibliography}

% Add references here, e.g.:
% \begin{thebibliography}{9}
% \bibitem{ref1} Author, Title, Publisher, Year.
% \end{thebibliography}

\end{document}
